\documentclass[11pt]{article}
\usepackage[a4paper,margin=25mm]{geometry}
\usepackage{amsmath,amssymb,amsthm,booktabs,array,longtable}
\usepackage[hidelinks]{hyperref}
\newtheorem{theorem}{Theorem}
\newtheorem{lemma}[theorem]{Lemma}
\newtheorem{corollary}[theorem]{Corollary}
\newcommand{\floor}[1]{\left\lfloor #1\right\rfloor}
\newcommand{\ceil}[1]{\left\lceil #1\right\rceil}
\title{Oudra's cubic and quadratic floor conjectures\\
and an alternative proof of Wu's tetrahedral formula}
\author{Haobo Ma\\\small ChronoAI Pte Ltd, Singapore
\and Wenlin Zhang\\\small National University of Singapore\\\small ChronoAI Pte Ltd, Singapore}
\date{November 14, 2026}
\begin{document}
\maketitle
\begin{abstract}
We prove Ridouane Oudra's cubic floor conjecture, posed on October 6,
2023 in the FORMULA field of OEIS A248186, revision 18, and give an
alternative proof of the tetrahedral counting formula contributed by
Chai Wah Wu to A056556 on November 4, 2024 and treated in his 2025
article, Section 4, p.~22. The least positive index $k$ for which
$1/18-\sum_{h=1}^k[h(h+1)(h+2)(h+3)]^{-1}<n^{-3}$ is
$\floor{rn+r^2/n-1}$ for $n>2$, where $r=3^{-1/3}$.
An integer-target cubic inversion also recovers Oudra's earlier
tetrahedral formula. A congruence argument proves his quadratic
conjecture in A248183. For both threshold sequences we obtain exact
jump positions and complementary densities, and for the cubic jumps
we characterize their deviations from the associated Beatty sequence.
\end{abstract}

\section{An integer gap and Oudra's cubic conjecture}
OEIS A248186 was introduced by Clark Kimberling in 2014; its floor
conjecture is credited to Ridouane Oudra \cite{cubic}. For $n\geq1$, set
\[
 R(k)=\frac1{18}-\sum_{h=1}^k\frac1{h(h+1)(h+2)(h+3)},\qquad
 a(n)=\min\{k\geq1:R(k)<n^{-3}\}.
\]
All indices are integers. The minimum is over positive $k$, as in the
entry's program. The following inversion works for every positive
integer target, rather than just for perfect cubes.

\begin{lemma}\label{lem:integer}
Let $Q(x)=3(x^3-x)$. For each integer $M\geq1$, put
\[
 t=(M/3)^{1/3},\qquad H(M)=\floor{t+\frac1{3t}}.
\]
Then $H(M)$ is the largest positive integer $m$ with $Q(m)\leq M$;
equivalently,
\begin{equation}\label{eq:bracket}
 Q(H(M))\leq M<Q(H(M)+1).
\end{equation}
\end{lemma}
\begin{proof}
For $y>x\geq1$,
\[
 Q(y)-Q(x)=3(y-x)(y^2+xy+x^2-1)>0.
\]
Write $b=t+1/(3t)$ and $m=\floor b$. Since
$b\geq2/\sqrt3>1$, we have $m\geq1$. Expansion gives
\begin{equation}\label{eq:overshoot}
 Q(b)=3t^3+\frac1{9t^3}=M+\frac1{3M}.
\end{equation}
As $m\leq b<m+1$, monotonicity yields
$Q(m)\leq M+1/(3M)<M+1$ and $M<Q(m+1)$.
The integers $Q(m)$ and $M$ therefore satisfy \eqref{eq:bracket}.
Strict monotonicity proves the maximality assertion.
\end{proof}

\begin{theorem}\label{thm:cubic}
The minimum defining $a(n)$ exists for every positive integer $n$.
We have $a(1)=a(2)=1$ and, with $r=3^{-1/3}$,
\begin{equation}\label{eq:cubic}
 a(n)=\floor{rn+\frac{r^2}{n}-1}\qquad(n\geq3).
\end{equation}
More generally, at a reciprocal integer target $1/M$, the least
positive $k$ with $R(k)<1/M$ is $\max\{1,H(M)-1\}$.
\end{theorem}
\begin{proof}
The identity
\[
 \frac1{h(h+1)(h+2)(h+3)}
 =\frac13\left(\frac1{h(h+1)(h+2)}-
                   \frac1{(h+1)(h+2)(h+3)}\right)
\]
telescopes, including at $k=0$, to
\begin{equation}\label{eq:tail}
 R(k)=\frac1{3(k+1)(k+2)(k+3)}.
\end{equation}
Thus $R(k)\to0$ and $R(k)<1/M$ is equivalent to $M<Q(k+2)$.
Lemma~\ref{lem:integer} makes the least positive admissible $k$ equal
to $\max\{1,H(M)-1\}$. For $M=n^3$ and $n\geq3$, $Q(2)=18<n^3$,
so $H(n^3)\geq2$. In this case $t=rn$ and $1/(3t)=r^2/n$,
which gives \eqref{eq:cubic}. For $n=1,2$, $R(1)=1/72<n^{-3}$,
so the positive minimum is $1$.
\end{proof}

\section{Cubic plateaus, jumps, and the Beatty comparison}
The COMMENT field of A248186 asserts zero--one increments and a
partition into the plateau and jump positions, A248187 and A248188
\cite{cubic,plateaus,jumps}.

\begin{corollary}\label{cor:jumps}
For every $n\geq1$, $a(n+1)-a(n)\in\{0,1\}$. The plateau positions
$\{n:a(n+1)=a(n)\}$ and jump positions $\{n:a(n+1)=a(n)+1\}$
partition the positive integers and have natural densities $1-r$
and $r$, respectively. The $j$th jump position is
\begin{equation}\label{eq:jump}
 c_j=\ceil{\sqrt[3]{3(j+1)(j+2)(j+3)}}-1
     =\floor{\sqrt[3]{3(j+1)(j+2)(j+3)-1}},\qquad j\geq1.
\end{equation}
\end{corollary}
\begin{proof}
For $n\geq3$, $b(n)=rn+r^2/n$ satisfies
\[
 0<b(n+1)-b(n)=r-\frac{r^2}{n(n+1)}<r<1.
\]
Taking floors proves the increment assertion; the initial cases follow
from $a(1)=a(2)=a(3)=a(4)=1$. For $j,n\geq1$, \eqref{eq:tail} gives
\begin{equation}\label{eq:level}
 a(n)\leq j\quad\Longleftrightarrow\quad
 n^3<3(j+1)(j+2)(j+3).
\end{equation}
Since $a$ is unbounded with zero--one increments and starts at $1$,
the largest such $n$ is where the value changes from $j$ to $j+1$.
This proves the ceiling formula. For integer $L>0$, the largest integer
with cube $<L$ is $\floor{\sqrt[3]{L-1}}$, proving the other formula.
There are $a(N+1)-1$ jumps in $1,\ldots,N$; \eqref{eq:cubic} gives
$a(N+1)=r(N+1)+O(1)$. Division by $N$ proves both densities.
\end{proof}

A248188 records agreement with $\floor{3^{1/3}(j+2)}$ for
$1\leq j\leq655$ and an exception at $j=656$ \cite{jumps,beatty}.
It does not assert equality for all $j$.

\begin{corollary}\label{cor:beatty}
Let $x=j+2$, $u_j=\floor{3^{1/3}x}$ and $L=3(x^3-x)$.
For every $j\geq1$,
\begin{equation}\label{eq:correction}
 c_j=u_j-\boldsymbol{1}_{\{u_j^3\geq L\}},
\end{equation}
where $\boldsymbol{1}$ denotes the indicator of the stated condition.
Thus the discrepancy is always zero or one, with an exact integer test.
\end{corollary}
\begin{proof}
Put $v=3^{1/3}x$ and $w=L^{1/3}$. For $x\geq3$,
\[
 (v-1)^3<L<v^3,
\]
because $L-(v-1)^3=3v^2-3v-3x+1>0$: here $v>x$ and
$v(v-1)>x(x-1)\geq2x$. Hence $v-1<w<v$.
Since $u_j\leq v<u_j+1$ and $v$ is irrational, $w>u_j-1$.
If $u_j^3<L$, then $u_j<w<u_j+1$ and $c_j=\ceil w-1=u_j$.
Otherwise $u_j-1<w\leq u_j$ and $c_j=u_j-1$.
\end{proof}

The integer comparisons in the accompanying program confirm the
reported range and give $c_{656}=948$, whereas $u_{656}=949$.
The universal equality $c_j=u_j$ is not asserted in the source;
it fails at this example. The examples in A248187 and A248188 end
with a complement-name slip: the complementary pair is A248187 and
A248188. In fact, $4$ belongs both to the value set of A248186
($a(8)=4$) and to A248188 ($c_1=4$), so those two sets are not
complementary.

\section{The quadratic conjecture and both of its siblings}
A248183, cross-referenced by A248186, contains Oudra's conjecture
dated September 6, 2023, in its FORMULA field \cite{quadratic}.
Define, for $n\geq1$,
\[
 d(n)=\min\left\{k\geq1:\frac14-
 \sum_{h=1}^k\frac1{h(h+1)(h+2)}<n^{-2}\right\}.
\]
\begin{corollary}\label{cor:quadratic}
The minimum exists, $d(1)=1$, and
\begin{equation}\label{eq:quadratic}
 d(n)=\floor{\sqrt{n^2/2+1}-\frac12}\qquad(n\geq2).
\end{equation}
The plateau and jump positions, A248184 and A248185, partition the
positive integers with respective densities $1-1/\sqrt2$ and
$1/\sqrt2$. The $j$th jump position is
\begin{equation}\label{eq:quadraticjump}
 e_j=\ceil{\sqrt{2(j+1)(j+2)}}-1
     =\floor{\sqrt{2(j+1)(j+2)-1}},\qquad j\geq1.
\end{equation}
\end{corollary}
\begin{proof}
Telescoping gives residual $1/[2(k+1)(k+2)]$, tending to zero.
Admissibility is $n^2<P(k+1)$, where $P(x)=2x(x+1)$.
For $n\geq2$, set $b=\sqrt{n^2/2+1}-1/2>1$ and $m=\floor b$.
Then $P(b)=n^2+3/2$ and $P(m)\leq P(b)<P(m+1)$.
For integer $m$, $P(m)$ is divisible by $4$, while $n^2$ is
$0$ or $1$ modulo $4$. A positive difference $P(m)-n^2$ would
therefore be at least $3$, contradicting $P(m)\leq n^2+3/2$.
It follows that $P(m)\leq n^2<P(m+1)$, so the least admissible
positive $k$ is $m$. At $n=1$ the residual for $k=1$ is $1/12<1$.

For $f(x)=\sqrt{x^2/2+1}-1/2$ and $x>0$,
$0<f'(x)=x/(2\sqrt{x^2/2+1})<1/\sqrt2<1$.
Thus the floors in \eqref{eq:quadratic} have zero--one increments
for $n\geq2$; $d(1)=d(2)=d(3)=1$ handles the boundary.
For positive $j,n$,
\[
 d(n)\leq j\quad\Longleftrightarrow\quad n^2<2(j+1)(j+2).
\]
The largest such $n$ is the $j$th jump position, giving
\eqref{eq:quadraticjump} exactly as in \eqref{eq:jump}.
There are $d(N+1)-1$ jumps up to $N$, and
$d(N+1)=(N+1)/\sqrt2+O(1)$, proving the densities.
\end{proof}

The NAME fields of A248184 and A248185 are precisely these two
position sets \cite{quadraticpositions}. In A248184's COMMENT,
the final complement name should be A248184, as correctly stated in
A248185's example: the nonempty set A248185 is not its own complement.

\section{Earlier tetrahedral formulas and verification scope}
Let $T(N)$ be A056556: each integer $m\geq0$ occurs
$(m+1)(m+2)/2$ times, starting at $N=0$. Then $T(N)$ is the
largest $m$ with $\binom{m+2}{3}\leq N$. Oudra supplied a radical
formula in 2021 and the simpler expression below in 2022;
Wu supplied the integer-root alternative in 2024 \cite{tetrahedral}.
Wu's general repeated-term treatment appears in his preprint and
published article \cite[Section 4, p.~22]{wu}. These are prior
counting results, not proofs of the two threshold conjectures.

\begin{corollary}\label{cor:tetrahedral}
For $N\geq1$, with $s=(6N)^{1/3}$ and $z=\floor{(6N+6)^{1/3}}$,
\[
 T(N)=\floor{s+\frac1{3s}-1}
 =\begin{cases}
 z,&N\geq\binom{z+2}{3},\\
 z-1,&N<\binom{z+2}{3}.
 \end{cases}
\]
Also $T(0)=0$.
\end{corollary}
\begin{proof}
Writing $x=m+1$ turns $\binom{m+2}{3}\leq N$ into
$Q(x)\leq18N$. Lemma~\ref{lem:integer} therefore gives
$T(N)=H(18N)-1$, which is the first formula.
If $l=T(N)$, then
\[
 l^3\leq l(l+1)(l+2)\leq6N<6(N+1)
 \leq(l+1)(l+2)(l+3)<(l+2)^3.
\]
The penultimate inequality holds because the two bounding products
are multiples of $6$. Hence $l\leq z\leq l+1$; the test
$\binom{z+2}{3}\leq N$ distinguishes these two cases.
\end{proof}

A074279, in which $m$ occurs $m^2$ times, is another related
inverse-counting example: Oudra gave a reciprocal-correction formula
in 2023 and Wu proved an integer-root alternative
\cite{squareblocks,wu}. Those formulas are existing results rather
than additional threshold conjectures. The divisibility question
A248189, linked from A248187, concerns $mn\mid\sigma(m^2+n^2)$,
where $\sigma$ is the divisor-sum function. The reciprocal-product
argument gives no conclusion about that question \cite{unrelated}.

\begin{sloppypar}
The pinned Lean source is
\href{https://github.com/the-omega-institute/trureturing/blob/f1c20718fca66e9ef0a4cb5f4f6a9a55badd1ab5/D5/S3/ArithSums/OudraA248186CubicFloor.lean}{\texttt{OudraA248186CubicFloor.lean}}.
In namespace \texttt{D5.S3.ArithSums}, the declaration
\texttt{OudraA248186CubicFloor.}\allowbreak\texttt{result}
proves existence of a
natural-valued sequence of least positive admissible indices at every
positive input, and \eqref{eq:cubic} for every such sequence at $n>2$.
This covers the existence and cubic-formula assertions of
Theorem~\ref{thm:cubic}, not its general integer-target extension.
Lemma~\ref{lem:integer} and Corollaries~\ref{cor:jumps}--\ref{cor:tetrahedral}
have the written proofs above. The accompanying standard-library
Python program supplies separate exact finite checks of the formulas,
inversion brackets, jump criteria and quoted OEIS comparisons;
its output specifies the ranges. These finite checks are distinct
from the statements for unbounded indices.
\end{sloppypar}

\small
\begin{longtable}{@{}>{\raggedright\arraybackslash}p{.28\linewidth}p{.65\linewidth}@{}}
\toprule
Source item & Coverage \\
\midrule
\endhead
A248186, FORMULA & Proved: Oudra's conjecture, Theorem~\ref{thm:cubic}. \\
Integer-target extension & Not posed; proved: Lemma~\ref{lem:integer} and the general reciprocal threshold in Theorem~\ref{thm:cubic}. \\
A248186, COMMENT & Proved: zero--one increments and partition, Corollary~\ref{cor:jumps}. \\
A248187/188, NAME & Derived: exact jumps and densities, Corollary~\ref{cor:jumps}; plateaus are their complement. \\
A248187/188, example & Refuted-and-corrected: complement-name assertion; overlap at $4$, corrected pair stated above. \\
A248188, COMMENT & Confirmed on the stated finite range $1$--$655$ and at $656$; the universal Beatty equality is not asserted in the source and fails at $656$. \\
A248188/A059539 & Not posed; proved: universal zero--one discrepancy test, Corollary~\ref{cor:beatty}. \\
A248183, FORMULA & Proved: Oudra's conjecture, Corollary~\ref{cor:quadratic}. \\
A248183, COMMENT; A248184/185,\newline NAME & Proved: quadratic partition; derived jump formula and densities, Corollary~\ref{cor:quadratic}. \\
A248184, COMMENT & Refuted-and-corrected: the nonempty jump set is not its own complement; corrected pair stated above. \\
A056556, FORMULA & Reproved: Oudra's 2022 formula and Wu's 2024 integer-root alternative; related published treatment Wu (2025), Corollary~\ref{cor:tetrahedral}. \\
A074279, FORMULA & Prior proof cited: Wu (preprint, 2024; journal, 2025), Section 4; no additional claim here. \\
A248189, COMMENT & Not covered: divisibility existence conjecture. \\
\bottomrule
\end{longtable}
\normalsize

\begin{thebibliography}{10}
\bibitem{cubic} C.~Kimberling, OEIS A248186, revision 18
(October 29, 2023), NAME, COMMENT and FORMULA; conjecture by
R.~Oudra (October 6, 2023).
\url{https://oeis.org/A248186/internal}.
\bibitem{plateaus} C.~Kimberling, OEIS A248187, revision 4
(October 6, 2014), NAME and example. \url{https://oeis.org/A248187}.
\bibitem{jumps} C.~Kimberling, OEIS A248188, revision 5
(October 6, 2014), NAME and COMMENT. \url{https://oeis.org/A248188}.
\bibitem{beatty} OEIS A059539, revision 35 (May 30, 2026), NAME.
\url{https://oeis.org/A059539}.
\bibitem{quadratic} C.~Kimberling, OEIS A248183, revision 19
(November 11, 2024), FORMULA; conjecture by R.~Oudra
(September 6, 2023). \url{https://oeis.org/A248183}.
\bibitem{quadraticpositions} C.~Kimberling, OEIS A248184 and A248185,
revision 4 of each (October 6, 2014), NAME and examples.
\url{https://oeis.org/A248184}; \url{https://oeis.org/A248185}.
\bibitem{tetrahedral} H.~Bottomley, OEIS A056556, revision 52
(November 9, 2025), FORMULA; contributions by R.~Oudra
(2021, 2022) and C.~W.~Wu (2024). \url{https://oeis.org/A056556}.
\bibitem{wu} C.~W.~Wu, Algorithms for Complementary Sequences,
\emph{Integers} \textbf{25} (2025), A95, 26 pp., published November 5, 2025.
\url{https://doi.org/10.5281/zenodo.17535229}.
Preprint: \url{https://arxiv.org/abs/2409.05844}; see Section 4.
\bibitem{squareblocks} J.~Perry, OEIS A074279, revision 85
(November 7, 2025), FORMULA; contributions by R.~Oudra (2023)
and C.~W.~Wu (2024). \url{https://oeis.org/A074279}.
\bibitem{unrelated} Z.-W.~Sun, OEIS A248189, revision 7
(October 3, 2014), NAME and COMMENT. \url{https://oeis.org/A248189}.
\end{thebibliography}
\noindent\small\textcopyright\ 2026 Haobo Ma and Wenlin Zhang.
Licensed under Creative Commons Attribution 4.0 International (CC BY 4.0).
\end{document}
